\documentclass[10pt,letterpaper]{article}
\usepackage{cornell-notes}

\title{Clause 05.02.09: Architecture Description Ad Elements}
\author{}
\date{}
\setDocTitle{Clause 05.02.09: Architecture Description Ad Elements}
\setDocSubtitle{ISO/IEC/IEEE 42010:2022}
\setDocOwner{ISO 42010 Cornell Notes}
\setDocDate{\today}
\setCornellCollection{ISO/IEC/IEEE 42010:2022 Cornell Notes}
\setCornellUnitType{Clause}
\setCornellUnitNumber{05.02.09}
\setCornellUnitTitle{Architecture Description Ad Elements}
\hypersetup{pdftitle={Clause 05.02.09: Architecture Description Ad Elements},pdfsubject={Cornell notes},pdfauthor={}}

\begin{document}
\maketitle
\begin{CornellOverviewBox}{Overview}
{\Large\bfseries\textcolor{CNavy}{Section 5.2.9: Architecture description (AD) elements}}\\[3pt]
{\small\textbf{Classification:} Conceptual foundation}\\
{\small\textbf{Learning objective:} Explain the role of architecture description (ad) elements and connect it to related architecture-description concepts.}
\end{CornellOverviewBox}
\vspace{3pt}

\begin{CornellNotesTable}
\CornellNoteRow{Purpose / key concept}{An AD element is an occurrence of one or more architectural concepts in an AD.}
\CornellNoteRow{Core details}{\begin{itemize}[leftmargin=*,nosep,topsep=1pt]
\item The AD elements include occurrences of the following architectural concepts: stakeholder, concern, aspect, stakeholder perspective, architecture viewpoint, architecture view, model kind, legend, architecture view component, architecture decision, architecture rationale and any correspondence and correspondence method specified on those constructs.
\item Any one of these concepts can have multiple occurrences in one or more AD.
\item An AD element occurrence specifies one or more architectural concepts.
\item An AD element in an AD can be refined or elaborated by reference to another AD.
\end{itemize}}
\CornellNoteRow{Example / application}{AD elements introduced by viewpoints or model kinds include use case constructs such as preconditions, actors, boundaries, systems; activity model constructs such as activities, inputs, outputs, controls, and mechanisms; architecture or design patterns to be employed.}
\CornellNoteRow{Related sections}{5.2.7, 5.2.8}
\CornellNoteRow{Active recall}{\begin{itemize}[leftmargin=*,nosep,topsep=1pt]
\item What is the central role of Architecture description (AD) elements?
\item How does Architecture description (AD) elements connect to adjacent architecture-description concepts?
\end{itemize}}
\end{CornellNotesTable}


\begin{CornellSummaryBox}{Summary}
An AD element is an occurrence of one or more architectural concepts in an AD.
\end{CornellSummaryBox}

\end{document}
